\section{Estado del arte}

\subsection{Accidentes nucleares y evolución de la seguridad nuclear}

El desarrollo de la tecnología nuclear durante la primera mitad del siglo XX estuvo
orientado inicialmente hacia aplicaciones militares \cite{Todeschini}. A partir de
la década de 1950 se impulsó activamente su uso civil, particularmente en la
generación de electricidad mediante reactores de fisión nuclear, lo que condujo a
la proliferación de plantas nucleares en diversas regiones del mundo. Este proceso
de expansión tecnológica trajo consigo riesgos inherentes que se manifestaron de
forma crítica en dos accidentes de escala global: Chernóbil (1986) y Fukushima
Daiichi (2011), ambos clasificados con el nivel máximo de gravedad (7) en la Escala
Internacional de Eventos Nucleares (INES) \cite{IAEAGSG152022}.

La literatura especializada coincide en que ninguno de estos accidentes puede
atribuirse a una causa única. En Chernóbil, la confluencia de deficiencias en el
diseño del reactor RBMK, violaciones de procedimientos operativos y la ausencia
de una cultura de seguridad sólida determinó tanto la ocurrencia del accidente
como las demoras en la respuesta \cite{Shlyakhter, Wakeford}. En Fukushima, la
vulnerabilidad de la planta ante la concurrencia de un terremoto de magnitud 9.0
y el tsunami subsiguiente evidenció que los planes de contingencia existentes no
contemplaban escenarios de desastres naturales simultáneos de tal magnitud
\cite{Durovic}. Estudios posteriores señalan que deficiencias organizacionales
dentro de la empresa operadora TEPCO contribuyeron además a la subestimación
sistemática de los riesgos \cite{Suzuki2014}.

Desde una perspectiva histórica más amplia, el accidente de Three Mile Island
(1979), clasificado en nivel 5 de la escala INES, constituyó el primer evento
severo en un reactor de agua a presión occidental y desencadenó una revisión
profunda de los criterios de diseño, los procedimientos operativos y los marcos
regulatorios en varios países \cite{Walker2004}. Aunque sus consecuencias
radiológicas externas fueron limitadas en comparación con Chernóbil, el accidente
demostró que la interacción entre fallos técnicos y errores humanos podía generar
escenarios fuera del diseño de base, anticipando lecciones que no fueron
suficientemente internalizadas antes de los accidentes posteriores \cite{Walker2004}.

Estos tres eventos constituyen los casos de referencia fundamentales en la
literatura sobre gestión de accidentes nucleares severos y han sido determinantes
en la evolución de los marcos regulatorios y las estrategias de gestión de
residuos radiactivos a escala internacional.

\subsection{Clasificación y caracterización de residuos radiactivos postaccidente}

La gestión de residuos radiactivos generados en accidentes nucleares severos
difiere sustancialmente de la gestión en condiciones operativas normales, tanto
en volumen como en la heterogeneidad de los materiales producidos. La IAEA
establece un sistema de clasificación basado en la actividad específica de los
residuos y en la vida media de los radionúclidos que contienen, distinguiendo
cuatro categorías principales: residuos de muy baja actividad (VLLW), residuos
de baja actividad (LLW), residuos de media actividad (ILW) y residuos de alta
actividad (HLW) \cite{IAEAGSG12009}. Esta clasificación es fundamental para
determinar las estrategias de almacenamiento, tratamiento y disposición final
aplicables a cada flujo de residuos.

Los accidentes nucleares severos generan simultáneamente residuos de todas estas
categorías, incluyendo materiales que no están contemplados en los flujos
convencionales de gestión. Entre ellos se incluyen suelos contaminados superficiales
(clasificados generalmente como VLLW o LLW), equipos de protección personal,
estructuras de confinamiento parcialmente contaminadas, resinas de intercambio
iónico saturadas (ILW), materiales filtrantes y, en los casos más severos,
combustible nuclear dañado o fundido (HLW) \cite{IAEAGSG12009, Ojovan2011}.
La coexistencia de todos estos flujos en un mismo emplazamiento exige sistemas de
gestión integrados capaces de manejar materiales con características físicas,
químicas y radiológicas muy diversas.

Steinhauser et al. \cite{Steinhauser2014} realizaron una comparación sistemática
de los inventarios de radionúclidos liberados en Chernóbil y Fukushima, concluyendo
que, si bien el accidente de Chernóbil liberó una fracción mayor del inventario
total del núcleo del reactor, Fukushima presentó una mayor dispersión de
radionúclidos volátiles como el $^{137}$Cs y el $^{131}$I hacia la atmósfera y
el océano. Esta diferencia tiene implicaciones directas para la caracterización de
los residuos y la planificación de las estrategias de remediación, ya que
determina la distribución espacial de la contaminación y los medios ambientales
prioritarios de intervención.

\subsection{Gestión del combustible nuclear dañado}

Uno de los desafíos técnicos más complejos en la gestión postaccidente es el
manejo del combustible nuclear dañado. En accidentes con fusión total o parcial
del núcleo del reactor, los materiales combustibles, las estructuras metálicas
del núcleo y los productos de fisión se funden y mezclan formando una masa
heterogénea denominada corium, cuyas propiedades fisicoquímicas dependen de la
composición específica del reactor y de las condiciones de temperatura alcanzadas
durante el accidente \cite{Journeau2012}. El corium presenta una actividad
radiactiva extremadamente elevada, genera calor residual de forma continua durante
décadas y contiene radionúclidos de larga vida media que requieren aislamiento del
entorno durante períodos de miles de años \cite{Journeau2012}.

En Chernóbil, el corium formado durante la explosión del reactor se solidificó
parcialmente en las estructuras inferiores del edificio del reactor, formando
las denominadas ``lavas de combustible''. Los estudios sobre su composición y
estabilidad a largo plazo indican que estos materiales continúan siendo
radiológicamente significativos y representan uno de los principales desafíos
para las futuras operaciones de desmantelamiento \cite{Borovoi1990}. En Fukushima,
la localización exacta y el estado del corium en los tres reactores dañados
continúa siendo objeto de investigación mediante técnicas de muografía de
muones y sistemas robóticos de inspección, dado que el acceso directo a las
estructuras internas de los reactores permanece inviable por las condiciones
radiológicas \cite{IAEA2015Fukushima}.

La extracción del combustible dañado de los reactores de Fukushima constituye
la fase más compleja y prolongada del proceso de desmantelamiento, con un horizonte
temporal estimado de entre 30 y 40 años \cite{IAEA2015Fukushima}. La gestión
posterior del combustible extraído, incluyendo su caracterización, acondicionamiento
y disposición final como HLW, representa uno de los vacíos más relevantes en la
planificación actual de la gestión postaccidente de Fukushima.

\subsection{Gestión postaccidente en Chernóbil}

La respuesta al accidente de Chernóbil implicó la implementación de una de las
operaciones de confinamiento y gestión de residuos más complejas documentadas en
la historia de la ingeniería nuclear. En la fase inmediata, la construcción del
sarcófago de hormigón sobre el reactor destruido tuvo como objetivo principal
reducir la dispersión de radionúclidos al ambiente y proporcionar condiciones
mínimas de seguridad para las labores de limpieza. Sin embargo, análisis
posteriores determinaron que la estructura presentaba deficiencias de hermeticidad
que permitían la entrada de agua pluvial, acelerando la corrosión de las armaduras
y generando riesgo de colapso estructural \cite{EBRD2006}.

El deterioro del sarcófago motivó el desarrollo del Nuevo Confinamiento Seguro
(New Safe Confinement, NSC), una estructura de arco metálico de 108 metros de
altura, 257 metros de longitud y 36.000 toneladas de peso, completada en 2016 y
deslizada sobre el reactor existente sin necesidad de demoler previamente el
sarcófago original. El NSC fue diseñado para una vida útil mínima de 100 años
y tiene por objetivos contener el material radiactivo existente, proporcionar
condiciones ambientales controladas dentro del confinamiento, proteger de las
condiciones meteorológicas externas y albergar los sistemas de manipulación
remota necesarios para las futuras operaciones de caracterización y extracción
de residuos \cite{EBRD2006}. El proyecto fue financiado mediante el Fondo de
Refugio de Chernóbil, administrado por el Banco Europeo de Reconstrucción y
Desarrollo (BERD), con contribuciones de 45 países \cite{EBRD2006}.

Los programas de remediación del suelo implementados en la zona de exclusión de
30 km incluyeron la remoción mecánica de la capa superficial del suelo en áreas
habitadas y la construcción de instalaciones de almacenamiento temporal para los
residuos generados durante las operaciones de limpieza. La literatura documenta
que la concentración de $^{137}$Cs en los suelos de la zona de exclusión continúa
disminuyendo tanto por decaimiento radiactivo como por migración hacia capas más
profundas, aunque las tasas de reducción varían significativamente en función de
las características edáficas locales \cite{Konoplev}. Los informes de UNSCEAR
han concluido que las dosis de radiación recibidas por la mayoría de la población
afectada, con excepción de los trabajadores de emergencia y los niños expuestos
a $^{131}$I a través del consumo de leche, fueron relativamente bajas en
comparación con las proyecciones iniciales, aunque el impacto psicológico y
socioeconómico del accidente ha sido documentado como igualmente significativo
\cite{UNSCEAR2008}.

\subsection{Gestión postaccidente en Fukushima}

El accidente de Fukushima Daiichi generó desafíos específicos en materia de
gestión de residuos que no tenían precedente directo en la experiencia acumulada
tras Chernóbil. El principal problema fue la producción continua de grandes
volúmenes de agua contaminada, resultado del enfriamiento activo y permanente de
los reactores dañados, así como de la infiltración de agua subterránea en los
edificios de los reactores. Para abordar esta problemática, TEPCO implementó el
Advanced Liquid Processing System (ALPS), un sistema de tratamiento que combina
procesos de adsorción, coprecipitación, filtración y separación química para
reducir las concentraciones de más de sesenta radionúclidos hasta niveles
compatibles con los estándares regulatorios \cite{IAEAALPS2025}.

Una limitación técnica relevante del sistema es su incapacidad para eliminar el
tritio de forma eficaz, dado que este isótopo se encuentra incorporado en las
propias moléculas de agua en forma de agua tritiada (HTO) y no puede ser separado
mediante los métodos convencionales de tratamiento \cite{IAEAALPS2025}. En 2023,
tras la revisión técnica independiente realizada por la IAEA, que concluyó que el
plan cumplía con los estándares internacionales de seguridad, las autoridades
japonesas iniciaron el vertido controlado al océano Pacífico del agua tratada
mediante el sistema ALPS, generando un debate científico e internacional sobre
los efectos radiológicos a largo plazo en los ecosistemas marinos \cite{IAEA2023water}.
Los materiales adsorbentes y filtros saturados generados durante el proceso de
tratamiento constituyen residuos radiactivos sólidos de media actividad que
requieren almacenamiento especializado, añadiendo un flujo adicional de residuos
al inventario del emplazamiento \cite{IAEAALPS2025}.

La dispersión terrestre de radionúclidos requirió una respuesta de remediación
a gran escala. Las autoridades japonesas implementaron programas de remoción del
suelo superficial contaminado en las zonas habitadas próximas a la planta,
generando aproximadamente 14 millones de metros cúbicos de residuos de baja y
media actividad que fueron almacenados temporalmente en contenedores distribuidos
por la prefectura de Fukushima \cite{MOE2021}. La definición de un sitio de
almacenamiento definitivo para estos residuos, con un plazo legal de traslado
fuera de la prefectura establecido para 2045, constituye uno de los desafíos
pendientes más relevantes en la gestión postaccidente \cite{MOE2021}.

Los estudios sobre la dispersión ambiental de radionúclidos en Fukushima han
documentado que factores climáticos como los tifones y las precipitaciones
intensas han redistribuido la contaminación de manera distinta a la prevista por
los modelos teóricos de decaimiento \cite{Saito}. La proximidad de la planta al
océano Pacífico introdujo una dimensión de contaminación marina sin precedente
en Chernóbil, con implicaciones para los ecosistemas acuáticos documentadas en
estudios de bioacumulación de $^{137}$Cs en especies marinas \cite{Buesseler2012}.

\subsection{Estrategias de descontaminación y remediación ambiental}

La descontaminación de suelos, edificaciones y cuerpos de agua contaminados por
accidentes nucleares ha sido objeto de una extensa investigación. Entre las
estrategias documentadas en la literatura se distinguen métodos de remoción
física, métodos de inmovilización química y métodos biológicos, cuya aplicabilidad
depende del tipo de radionúclido, las características del substrato y la escala
espacial de la intervención requerida \cite{Vandenhove2009}.

La remoción mecánica de la capa superficial del suelo es la técnica más ampliamente
aplicada en accidentes nucleares, dado que la mayor parte de la contaminación
por radionúclidos como el $^{137}$Cs y el $^{90}$Sr se concentra en los primeros
centímetros del perfil edáfico \cite{Vandenhove2009}. Si bien esta técnica es
eficaz para reducir rápidamente las tasas de dosis en superficies habitadas,
genera grandes volúmenes de residuos de baja actividad que deben ser gestionados
en instalaciones adecuadas. En Fukushima, la aplicación sistemática de esta
técnica en zonas urbanas y agrícolas generó los volúmenes de residuos antes
mencionados \cite{MOE2021}.

La fitorremediación, que utiliza especies vegetales con capacidad de absorber
radionúclidos del suelo y acumularlos en su biomasa, ha sido investigada como
alternativa complementaria a los métodos físicos, especialmente para el tratamiento
de áreas extensas con niveles de contaminación moderados \cite{Dushenkov2010}.
Sin embargo, su eficacia está limitada por las tasas de absorción de cada especie,
las propiedades de retención del suelo y la necesidad de gestionar la biomasa
contaminada resultante como residuo radiactivo, lo que reduce su ventaja comparativa
en términos de generación de residuos secundarios.

\subsection{Almacenamiento temporal y disposición definitiva de residuos}

El almacenamiento de los residuos radiactivos generados por accidentes nucleares
plantea desafíos que se extienden más allá del horizonte temporal de la respuesta
inmediata. La distinción entre almacenamiento temporal, que implica la posibilidad
de recuperación futura de los residuos, y disposición definitiva, que supone su
aislamiento permanente del entorno humano, es fundamental en la planificación
estratégica de la gestión postaccidente \cite{IAEAWSR62021}.

Para los residuos de alta actividad y el combustible irradiado, el repositorio
geológico profundo (deep geological repository, DGR) es reconocido por la
comunidad científica internacional y por la IAEA como la solución técnicamente
más robusta disponible, al proporcionar el aislamiento necesario durante los
períodos de miles a millones de años que requieren estos materiales
\cite{Chapman2002}. Sin embargo, la implementación de repositorios geológicos
profundos enfrenta desafíos técnicos, regulatorios y sociales que han retrasado
su materialización en la mayoría de los países con programas nucleares activos
\cite{Chapman2002}. A la fecha de redacción de este trabajo, Finlandia y Suecia
son los únicos países que cuentan con emplazamientos seleccionados y en fase
avanzada de construcción para repositorios geológicos de HLW.

En el contexto específico de los accidentes nucleares, ninguno de los países
más afectados (Ucrania para Chernóbil y Japón para Fukushima) cuenta con
repositorios geológicos profundos operativos para la disposición final de los
residuos de alta actividad generados \cite{IAEAWSR62021}. Esta situación implica
que los residuos más peligrosos continúan en almacenamiento temporal en los
propios emplazamientos accidentados, con los riesgos a largo plazo que ello
conlleva.

\subsection{Cultura de seguridad y reformas regulatorias internacionales}

La literatura postaccidente ha consolidado el concepto de cultura de seguridad
como variable analítica fundamental en la evaluación del riesgo nuclear. La IAEA
definió por primera vez este concepto en el informe INSAG-4 (1991) como el
conjunto de características y actitudes que, en organizaciones e individuos,
garantiza que los temas de seguridad nuclear reciban la atención que su
importancia merece \cite{INSAG4}. Posteriormente, el informe INSAG-15 amplió
este marco conceptual incorporando indicadores de desempeño y mecanismos de
autoevaluación para las organizaciones nucleares \cite{INSAG15}.

En el caso de Fukushima, estudios posteriores al accidente documentaron que la
cultura organizacional de TEPCO presentaba deficiencias que contribuyeron a la
subestimación sistemática de los riesgos y a la incapacidad de gestionar
adecuadamente escenarios de emergencia complejos \cite{Suzuki2014}. Este análisis
es consistente con las conclusiones sobre Chernóbil, donde la opacidad
institucional y la subordinación de los criterios técnicos a las presiones
políticas fueron identificadas como factores determinantes tanto en la ocurrencia
del accidente como en la gestión de sus consecuencias \cite{Shlyakhter, Wakeford}.

El accidente de Chernóbil impulsó la adopción de la Convención sobre Seguridad
Nuclear (1994) y la Convención Conjunta sobre Seguridad en la Gestión del
Combustible Gastado y sobre Seguridad en la Gestión de Desechos Radiactivos (1997),
instrumentos que establecieron obligaciones internacionales vinculantes en materia
de seguridad nuclear y gestión de residuos \cite{IAEAJointConvention}. Tras
Fukushima, la IAEA aprobó el Plan de Acción sobre Seguridad Nuclear (2011), que
identificó doce áreas de mejora prioritarias, incluyendo la evaluación de la
seguridad ante eventos externos extremos, el fortalecimiento de la gestión de
emergencias y la mejora de la comunicación de riesgos a la población
\cite{IAEAActionPlan2011}. La revisión de los requisitos de preparación y respuesta
ante emergencias nucleares derivó en la publicación de nuevos estándares de
seguridad que reforzaron los mecanismos de coordinación interinstitucional y los
criterios de diseño ante eventos que superen las previsiones iniciales
\cite{IAEAGSG152022}.

\subsection{Desafíos actuales y perspectivas de investigación}

La revisión de la literatura sobre gestión de residuos radiactivos en accidentes
nucleares revela un conjunto de desafíos que continúan siendo objeto de
investigación activa. En primer lugar, la ausencia de soluciones de almacenamiento
definitivo para los residuos de alta y media actividad generados en Chernóbil y
Fukushima representa uno de los problemas estructurales más relevantes, dado que
los repositorios geológicos profundos aún no han sido implementados en ninguno de
los dos países afectados \cite{IAEAWSR62021}.

En segundo lugar, la modelización del comportamiento a largo plazo de los
radionúclidos en entornos complejos continúa siendo un área de investigación
activa, dado que los modelos teóricos existentes han demostrado limitaciones para
predecir la redistribución de contaminantes en presencia de variables climáticas
y geológicas locales \cite{Saito, Konoplev}.

En tercer lugar, la literatura ha identificado una brecha significativa entre los
marcos normativos internacionales y las capacidades reales de implementación a
nivel nacional, especialmente en lo que respecta a la sostenibilidad financiera
e institucional de los programas de remediación a largo plazo \cite{IAEAGSG152022}.
Ojovan y Lee \cite{Ojovan2011} señalan que la gestión de residuos radiactivos en
contextos postaccidente requiere enfoques adaptativos que integren la incertidumbre
inherente sobre la evolución del inventario de residuos, las condiciones del
emplazamiento y los marcos regulatorios a lo largo de horizontes temporales de
décadas. Esta observación subraya la relevancia de desarrollar investigación
comparativa sobre gobernanza y financiamiento de la gestión postaccidente
aplicable a diferentes contextos institucionales y socioeconómicos.

La experiencia acumulada en Chernóbil y Fukushima ha generado un corpus
bibliográfico extenso, pero la literatura señala que el conocimiento adquirido
no ha sido transferido de manera sistemática hacia la planificación preventiva
en otras instalaciones nucleares activas \cite{Wakeford}. Esta observación
justifica la necesidad de continuar desarrollando investigación sobre la gestión
de residuos en accidentes nucleares, con el objetivo de fortalecer tanto los
marcos técnicos como los institucionales de respuesta ante futuros eventos de
esta naturaleza.
